\chapter{System Design}
All diagrams in this chapter were drawn from the structure of the code in the repository.

\section{Architecture}
\begin{figure}[H]
\centering
\resizebox{\textwidth}{!}{\begin{tikzpicture}
  % Presentation tier
  \node[box, minimum width=146mm, text width=140mm] (spa) {\textbf{React 19 single-page application} (Vite, TypeScript, Tailwind)\\[1pt]
    \footnotesize landing, sign-in, dashboard, portfolio, new analysis, result page, history, settings};

  % Application tier
  \node[box, below=14mm of spa.south west, anchor=north west, xshift=3mm, text width=30mm, minimum height=19mm] (api) {\textbf{api/v1}\\ \footnotesize auth, users, portfolio, decisions, market, admin};
  \node[box, right=3mm of api, text width=30mm, minimum height=19mm] (mw) {\textbf{middleware}\\ \footnotesize CSRF, rate limits, headers, logging};
  \node[box, right=3mm of mw, text width=30mm, minimum height=19mm] (auth) {\textbf{auth}\\ \footnotesize RS256 JWT, bcrypt, TOTP, Google OIDC};
  \node[box, right=3mm of auth, text width=30mm, minimum height=19mm] (svc) {\textbf{services}\\ \footnotesize portfolio, prices, market, e-mail, deletion};
  \node[box, below=3mm of api.south west, anchor=north west, text width=135mm] (pipe)
    {\textbf{core/pipeline.py} (background task): bias $\rightarrow$ simulation $\rightarrow$ explanation $\rightarrow$ recommendation $\rightarrow$ audit};
  \begin{scope}[on background layer]
    \node[tier, fit=(api)(svc)(pipe), inner sep=3mm] (apptier) {};
  \end{scope}
  \node[tierlabel, anchor=south west] at (apptier.north west) {Application tier: FastAPI (backend/)};

  % Analytics tier (left)
  \node[mlbox, below=15mm of apptier.south west, anchor=north west, xshift=3mm, text width=74mm] (src)
    {\textbf{Research modules (src/)}\\ \footnotesize dataset\_v4, bias\_model\_v4 (XGBoost, SHAP),\\ \footnotesize simulator\_v2, causal\_v4 (DoWhy),\\ \footnotesize explainer (English and Hindi)};
  \node[store, below=4mm of src, text width=74mm] (art)
    {\textbf{Model artefacts and result files}\\ \footnotesize trained model, causal results,\\ \footnotesize cached market history};
  \draw[flow] (art) -- (src);
  \begin{scope}[on background layer]
    \node[tier, fit=(src)(art), inner sep=3mm] (mltier) {};
  \end{scope}
  \node[tierlabel, anchor=south west] at (mltier.north west) {Analytics tier, via core/ml\_bridge.py};

  % Data tier (right)
  \node[store, anchor=north east, text width=50mm] (pg) at ($(apptier.south east)+(-3mm,-15mm)$)
    {\textbf{PostgreSQL 18}\\ \footnotesize row-level security};
  \node[store, below=3mm of pg, text width=50mm] (redis)
    {\textbf{Redis / Memurai}\\ \footnotesize sessions, rate limits, events};
  \node[ext, below=3mm of redis, text width=50mm] (ext)
    {\textbf{External services}\\ \footnotesize Yahoo Finance, MFAPI.in, Google;\\ \footnotesize Gemini and e-mail (optional)};
  \begin{scope}[on background layer]
    \node[tier, fit=(pg)(ext), inner sep=3mm] (datatier) {};
  \end{scope}
  \node[tierlabel, anchor=south west] at (datatier.north west) {Data and external services};

  \draw[biflow] (spa.south) -- node[lbl, right=1mm] {HTTPS, JSON, httpOnly cookies (/v1/*)} (spa.south |- apptier.north);
  \draw[biflow] ($(mltier.north east)+(-14mm,0)$) coordinate (m) -- node[lbl, left=1mm] {in-process calls} (m |- apptier.south);
  \draw[biflow] ($(datatier.north east)+(-10mm,0)$) coordinate (d) -- (d |- apptier.south);
\end{tikzpicture}
}
\caption{System architecture of Folium}
\label{fig:arch}
\end{figure}

Figure \ref{fig:arch} shows the overall architecture. The presentation tier is a React single-page application. During development the Vite server forwards every \texttt{/v1/*} request to the API, so the browser sees a single origin, which keeps the httpOnly session cookies and the CSRF check simple.

The application tier is a FastAPI service. Its routes are in \texttt{backend/api/v1/}, its cross-cutting checks in \texttt{backend/middleware/}, its authentication code in \texttt{backend/auth/} and its domain services in \texttt{backend/services/}. The analysis itself runs as a background task under \texttt{backend/core/pipeline.py}.

The analytics tier is the research code in \texttt{src/}. An important design decision is that the back end never re-implements the models. The module \texttt{backend/core/ml\_bridge.py} is the only place where the back end reaches into \texttt{src/}. It imports the heavy ML libraries lazily (they take about two seconds and 400 MB to load), runs blocking work in a worker thread so that the server stays responsive, and maps the vocabulary between the two sides. Since the research modules produced the reported numbers, re-implementing them would risk silently changing those numbers.

PostgreSQL stores users, portfolios, decisions and the audit chain under row-level security. A Redis-compatible store holds the token blacklist, the rate-limit counters and the event backlog for live progress. External services are called only for prices and market data, for Google sign-in, and optionally for e-mail and Gemini.

\section{Analysis Pipeline Components}
\begin{figure}[H]
\centering
\resizebox{0.95\textwidth}{!}{\begin{tikzpicture}[node distance=3mm]
  \node[box, text width=32mm] (dec) {\textbf{decisions.py}\\ \footnotesize POST /analyze (202)\\ \footnotesize GET /stream (SSE)};
  \node[box, below=10mm of dec, text width=32mm] (pipe) {\textbf{pipeline.py}\\ \footnotesize orchestrator};
  \node[box, below=10mm of pipe, text width=32mm] (ev) {\textbf{events.py}\\ \footnotesize Redis pub/sub\\ \footnotesize and backlog};
  \draw[flow] (dec) -- node[lbl, right] {schedules} (pipe);
  \draw[flow] (pipe) -- node[lbl, right] {publishes} (ev);

  \node[box, right=22mm of dec.north east, anchor=north west, text width=34mm, yshift=6mm] (s1) {\footnotesize 1. bias\_detector.py\\ \footnotesize + comment\_reader.py};
  \node[box, below=of s1, text width=34mm] (s2) {\footnotesize 2. simulator.py};
  \node[box, below=of s2, text width=34mm] (s3) {\footnotesize 3. explainer.py};
  \node[box, below=of s3, text width=34mm] (s4) {\footnotesize 4. recommendation.py\\ \footnotesize + debate.py};
  \node[box, below=of s4, text width=34mm] (s5) {\footnotesize 5. audit.py};
  \foreach \s in {s1,s2,s3,s4,s5} \draw[flow, greyline] (pipe.east) -- (\s.west);

  \node[mlbox, right=10mm of s2.east, anchor=west, text width=20mm, minimum height=32mm] (br) {\textbf{ml\_bridge}\\ \footnotesize lazy import,\\ \footnotesize worker thread,\\ \footnotesize label mapping};
  \foreach \s in {s1,s2,s3} \draw[flow, blueline] (\s.east) -- (br.west |- \s.east);

  \node[mlbox, right=8mm of br.north east, anchor=north west, text width=22mm] (m1) {\footnotesize bias\_model\_v4};
  \node[mlbox, below=2mm of m1, text width=22mm] (m2) {\footnotesize simulator\_v2};
  \node[mlbox, below=2mm of m2, text width=22mm] (m3) {\footnotesize causal\_v4};
  \node[mlbox, below=2mm of m3, text width=22mm] (m4) {\footnotesize explainer,\\ \footnotesize llm\_polish};
  \foreach \m in {m1,m2,m3,m4} \draw[flow, blueline] (br.east) -- (\m.west);
  \node[tierlabel, above=1mm of m1] {src/};

  \node[store, below=6mm of s5, text width=34mm] (al) {\footnotesize audit\_logs (PostgreSQL)};
  \draw[flow] (s5) -- (al);
\end{tikzpicture}
}
\caption{Components involved in one analysis}
\label{fig:components}
\end{figure}

Figure \ref{fig:components} shows the components of one analysis. The analyze route (\texttt{POST} on \texttt{/v1/decisions/analyze}) schedules the pipeline as a background task and returns immediately. The orchestrator calls five steps in a fixed order: bias detection, simulation, explanation, recommendation and audit. Each step returns the columns it wants to update and the fields to attach to its progress event, while the orchestrator alone owns the status, the events and persistence. It refuses any step that tries to write the columns it reserves for itself. The first three steps call the research code through the bridge, and the audit step writes to the \texttt{audit\_logs} table.

\section{Data-Flow Diagrams}
\begin{figure}[H]
\centering
\resizebox{0.9\textwidth}{!}{\begin{tikzpicture}
  \node[proc, minimum size=34mm, font=\small] (sys) {\textbf{0}\\ Folium\\ decision-support\\ system};
  \node[entity, left=34mm of sys] (inv) {Investor};
  \node[entity, right=30mm of sys, yshift=14mm] (mkt) {Market data\\ providers};
  \node[entity, right=30mm of sys, yshift=-14mm] (gem) {Gemini API\\ (optional)};
  \draw[flow] ([yshift=4mm]inv.east) -- node[lbl, above=1mm, align=center] {profile, holdings,\\ action, comment} ([yshift=4mm]inv.east -| sys.west);
  \draw[flow] ([yshift=-4mm]inv.east -| sys.west) -- node[lbl, below=1mm, align=center] {verdict, projections,\\ explanation, record} ([yshift=-4mm]inv.east);
  \draw[flow] (mkt.west) -- node[lbl, above, sloped] {prices, NIFTY, VIX} (sys.north east);
  \draw[flow, dashed] ([yshift=2mm]sys.south east) -- node[lbl, above, sloped] {comment, numbers} ([yshift=2mm]gem.west);
  \draw[flow, dashed] ([yshift=-2mm]gem.west) -- node[lbl, below, sloped] {checked wording} ([yshift=-2mm]sys.south east);
\end{tikzpicture}
}
\caption{Data-flow diagram, level 0}
\label{fig:dfd0}
\end{figure}

At the context level (Figure \ref{fig:dfd0}) the system exchanges data with three external entities. The investor supplies a profile, holdings, the intended action and an optional comment, and receives a verdict, projections, an explanation and a verifiable record. Market data providers supply prices, the NIFTY index and the VIX. The Gemini API is optional: it receives a comment and the analysis numbers only when the investor has consented and the server has enabled it, and its output is accepted only after checks.

\begin{figure}[H]
\centering
\resizebox{0.95\textwidth}{!}{\begin{tikzpicture}[proc/.append style={minimum size=24mm, text width=20mm}]
  \node[proc] (p1) {1.0\\ Capture\\ context};
  \node[proc, right=14mm of p1] (p2) {2.0\\ Read comment,\\ detect bias};
  \node[proc, right=14mm of p2] (p3) {3.0\\ Simulate\\ hold vs sell};
  \node[proc, below=16mm of p3] (p4) {4.0\\ Explain\\ (EN / HI)};
  \node[proc, left=14mm of p4] (p5) {5.0\\ Recommend,\\ both sides};
  \node[proc, left=14mm of p5] (p6) {6.0\\ Seal audit\\ record};
  \node[entity, left=10mm of p1] (inv) {Investor};
  \draw[flow] (inv) -- node[lbl, above] {request} (p1);
  \draw[flow] (p1) -- node[lbl, above] {snapshot} (p2);
  \draw[flow] (p2) -- node[lbl, above] {pattern} (p3);
  \draw[flow] (p3) -- node[lbl, right] {curves, $p$} (p4);
  \draw[flow] (p4) -- node[lbl, above] {narrative} (p5);
  \draw[flow] (p5) -- node[lbl, above] {verdict} (p6);
  \node[datastore, below=12mm of p6] (d3) {D3 \; audit\_logs};
  \node[datastore, below=12mm of p5] (d2) {D2 \; decisions};
  \node[datastore, below=12mm of p4] (d1) {D1 \; users, profiles};
  \draw[flow] (p6) -- node[lbl, right] {append} (d3);
  \draw[flow] (p5) -- node[lbl, right] {write outputs} (d2);
  \draw[flow] (d1.east) -- ++(6mm,0) |- node[lbl, pos=0.25, right] {read profile} (p3.east);
\end{tikzpicture}
}
\caption{Data-flow diagram, level 1 (the analysis process)}
\label{fig:dfd1}
\end{figure}

Figure \ref{fig:dfd1} breaks the analysis into six processes. Process 1.0 is the analyze endpoint, which captures the context, and processes 2.0 to 6.0 are the five pipeline steps. Each process writes its outputs into the decisions store, and the next process reads the refreshed row, which is why the orchestrator reloads the decision before every step. The simulation reads the profile for income, SIP and the destination of sale proceeds, and the last process appends an entry to the audit store.

\section{Use Case Diagram}
\begin{figure}[H]
\centering
\resizebox{0.95\textwidth}{!}{\begin{tikzpicture}
  \foreach \t [count=\i] in {Register or log in (e-mail or Google), Set up two-step verification, Maintain financial profile,
      Manage portfolio holdings, Request decision analysis, Watch live progress, Read explanation (English or Hindi),
      Download report (HTML or PDF), Verify audit record, Export data or delete account, Reset a user's two-step verification}
    \node[usecase] (u\i) at (0,-\i*0.95) {\t};
  \begin{scope}[on background layer]
    \node[draw=greyline, rounded corners=4pt, fit=(u1)(u11), inner sep=5mm] (sysb) {};
  \end{scope}
  \node[tierlabel, anchor=south] at (sysb.north) {Folium};
  \actor{inv}{-6.4,-5.0}{Investor}
  \actor{adm}{6.4,-10.0}{Administrator}
  \actor{mds}{6.4,-3.8}{Market data\\ service}
  \foreach \i in {1,...,10} \draw[greyline] (-6.1,-4.45) -- (u\i.west);
  \draw[greyline] (6.1,-9.45) -- (u11.east);
  \draw[greyline] (6.1,-3.25) -- (u4.east);
  \draw[greyline] (6.1,-3.25) -- (u5.east);
\end{tikzpicture}
}
\caption{Use case diagram}
\label{fig:usecase}
\end{figure}

Figure \ref{fig:usecase} shows the actors and their use cases. The investor can register or log in by e-mail or with Google and then uses every feature of the application. The administrator's implemented use case is resetting a user's two-step verification. The market data service takes part in portfolio management, through live prices, and in analysis requests, through the market context and the stock's measured risk.

\section{Activity Diagram}
\begin{figure}[H]
\centering
\resizebox{0.85\textwidth}{!}{\begin{tikzpicture}[node distance=5mm]
  \node[circle, fill=black, minimum size=4mm] (start) {};
  \node[act, below=of start] (a1) {Investor picks a holding, an action, months held and an optional reason};
  \node[act, below=of a1] (a2) {POST /v1/decisions/analyze: validate input, capture market,\\ portfolio snapshot and the stock's beta and own volatility};
  \node[act, below=of a2] (a3) {Insert the decision with status \emph{pending}; return 202 with its id};
  \node[act, below=of a3, fill=bluefill, draw=blueline] (b1) {Step 1: read the comment, then run bias model v4.1};
  \node[act, below=of b1, fill=bluefill, draw=blueline] (b2) {Step 2: simulate HOLD and SELL on shared market paths};
  \node[act, below=of b2, fill=bluefill, draw=blueline] (b3) {Step 3: cause and effect, then English and Hindi explanation};
  \node[act, below=of b3, fill=bluefill, draw=blueline] (b4) {Step 4: evidence rules give the verdict; build both sides};
  \node[act, below=of b4, fill=bluefill, draw=blueline] (b5) {Step 5: append the hash-chained audit entry};
  \node[act, below=of b5] (c) {Mark the decision \emph{complete}; publish the final event};
  \node[circle, draw=black, line width=1pt, minimum size=5mm, below=of c] (end) {};
  \node[circle, fill=black, minimum size=2.6mm] at (end) {};
  \foreach \x/\y in {start/a1,a1/a2,a2/a3,a3/b1,b1/b2,b2/b3,b3/b4,b4/b5,b5/c,c/end} \draw[flow] (\x) -- (\y);
  \node[box, fill=sandfill, draw=sandline, right=8mm of b3, text width=24mm, font=\footnotesize] (f) {Any step raises an error: status \emph{failed}, a \emph{failed} event is published};
  \begin{scope}[on background layer]
    \node[fit=(b1)(b5), inner sep=2mm] (bg) {};
  \end{scope}
  \draw[flow, dashed, sandline] (bg.east) -- (f.west);
\end{tikzpicture}
}
\caption{Activity diagram of a decision analysis}
\label{fig:activity}
\end{figure}

Figure \ref{fig:activity} traces one analysis. The synchronous part ends with an HTTP 202 response, so the browser never waits for the models. The five shaded steps run in the background. If any step raises an exception, the decision is marked as failed, the error is stored and a failed event is published so that the browser stops waiting. Problems with wording never fail an analysis, because the explanation step falls back to the template text.

\section{Sequence Diagram}
\begin{figure}[H]
\centering
\resizebox{\textwidth}{!}{\begin{tikzpicture}
  \foreach \n/\t [count=\i] in {b/Browser, api/{FastAPI\\ decisions.py}, p/{Pipeline\\ (background)}, r/Redis, db/PostgreSQL, m/{src/ models}}
    \node[box, text width=20mm, minimum height=11mm, font=\footnotesize] (\n) at (\i*2.5,0) {\t};
  \foreach \n in {b,api,p,r,db,m} \draw[greyline, dashed] (\n.south) -- ++(0,-12.4);
  \draw[flow, black] (b.south |- 0,-1.08) -- node[lbl, above, font=\scriptsize] {POST /analyze} (api.south |- 0,-1.08);
  \draw[flow, black] (api.south |- 0,-1.86) -- node[lbl, above, font=\scriptsize] {INSERT decision} (db.south |- 0,-1.86);
  \draw[flow, tealline] (api.south |- 0,-2.64) -- node[lbl, above, font=\scriptsize] {202, decision id} (b.south |- 0,-2.64);
  \draw[flow, black] (b.south |- 0,-3.42) -- node[lbl, above, font=\scriptsize] {GET /stream (SSE)} (api.south |- 0,-3.42);
  \draw[flow, black] (api.south |- 0,-4.20) -- node[lbl, above, font=\scriptsize] {replay backlog, subscribe} (r.south |- 0,-4.20);
  \draw[flow, black] (api.south |- 0,-4.98) -- node[lbl, above, font=\scriptsize] {start task} (p.south |- 0,-4.98);
  \draw[flow, black] (p.south |- 0,-5.76) -- node[lbl, above, font=\scriptsize] {SELECT decision (RLS)} (db.south |- 0,-5.76);
  \draw[flow, black] (p.south |- 0,-6.54) -- node[lbl, above, font=\scriptsize] {predict, simulate, explain} (m.south |- 0,-6.54);
  \draw[flow, black] (p.south |- 0,-7.32) -- node[lbl, above, font=\scriptsize] {UPDATE step columns} (db.south |- 0,-7.32);
  \draw[flow, black] (p.south |- 0,-8.10) -- node[lbl, above, font=\scriptsize] {publish step done} (r.south |- 0,-8.10);
  \draw[flow, tealline] (r.south |- 0,-8.88) -- node[lbl, above, font=\scriptsize] {event} (api.south |- 0,-8.88);
  \draw[flow, tealline] (api.south |- 0,-9.66) -- node[lbl, above, font=\scriptsize] {SSE: step status} (b.south |- 0,-9.66);
  \draw[flow, black] (p.south |- 0,-10.44) -- node[lbl, above, font=\scriptsize] {INSERT audit\_logs} (db.south |- 0,-10.44);
  \draw[flow, black] (p.south |- 0,-11.22) -- node[lbl, above, font=\scriptsize] {publish complete} (r.south |- 0,-11.22);
  \draw[flow, black] (b.south |- 0,-12.00) -- node[lbl, above, font=\scriptsize] {GET /result} (api.south |- 0,-12.00);
  \node[draw=greyline, dashed, rounded corners=2pt, fit={(2.0,-5.31) (16.0,-9.81)}, inner sep=0pt] (loop) {};
  \node[tierlabel, anchor=north west, fill=white, inner sep=1pt] at (loop.north west) {loop: five steps};
\end{tikzpicture}
}
\caption{Sequence diagram of an analysis with live progress}
\label{fig:sequence}
\end{figure}

Figure \ref{fig:sequence} shows the order of messages. One detail is the backlog replay. If the browser only subscribed to Redis publish and subscribe, any event published before the subscription would be lost, and reloading the page during an analysis would show an empty history. Every event is therefore also appended to a Redis list with an expiry time. The stream endpoint first replays this list and then follows live messages, removing duplicates by a sequence number. If the stream cannot be opened at all, the result page falls back to polling.

\section{Entity-Relationship Diagram}
\begin{figure}[H]
\centering
\resizebox{0.95\textwidth}{!}{\begin{tikzpicture}[
  ent/.style={rectangle split, rectangle split parts=2, draw=black, line width=0.6pt, align=flush left,
              rectangle split part fill={tealfill, white}, font=\footnotesize\raggedright, inner sep=3pt, text width=45mm},
  rel/.style={draw=greyline, line width=0.6pt}]
  \node[ent] (users) at (0,0) {\textbf{users}\nodepart{two}PK user\_id\\ email (unique)\\ password\_hash, mfa\_secret\_enc\\ age, income, EMIs, rent, fees\\ risk tolerance, horizon\\ ai\_consent, purge\_after\\ created\_at};
  \node[ent] (port) at (-6.3,2.0) {\textbf{portfolios}\nodepart{two}PK holding\_id\\ FK user\_id\\ symbol, exchange\\ quantity, avg\_buy\_price\\ asset\_type};
  \node[ent] (fin) at (6.3,2.0) {\textbf{financial\_profiles}\nodepart{two}PK, FK user\_id\\ gold, FDs, savings\\ property, loans\\ rent received, monthly SIP};
  \node[ent] (dec) at (-6.3,-4.6) {\textbf{decisions}\nodepart{two}PK decision\_id\\ FK user\_id\\ symbol, prices, intent, reason\\ market snapshot\\ portfolio\_context (jsonb)\\ comment\_facts (jsonb)\\ bias, confidence, curves\\ recommendation, status};
  \node[ent] (aud) at (6.3,-4.6) {\textbf{audit\_logs}\nodepart{two}PK log\_id\\ FK user\_id, FK decision\_id\\ payload\_hash\\ previous\_hash, chain\_hash\\ payload (jsonb)};
  \node[ent] (tok) at (0,-5.2) {\textbf{refresh\_tokens, mfa\_*}\nodepart{two}FK user\_id\\ token\_hash, jti\\ code\_hash, used};
  \node[ent] (log) at (0,4.1) {\textbf{profile\_change\_log}\nodepart{two}FK user\_id, field, old\_value\_hash};
  \draw[rel] (users.west) -- node[lbl] {1 : N} (port.east);
  \draw[rel] (users.east) -- node[lbl] {1 : 1} (fin.west);
  \draw[rel] (users.south west) -- node[lbl] {1 : N} (dec.north east);
  \draw[rel] (users.south east) -- node[lbl] {1 : N} (aud.north west);
  \draw[rel] (users.south) -- node[lbl] {1 : N} (tok.north);
  \draw[rel] (users.north) -- node[lbl] {1 : N} (log.south);
  \draw[rel] (dec.south) -- ++(0,-0.5) coordinate (q) -- node[lbl] {1 : N (one decision, its audit entries)} (q -| aud.south) -- (aud.south);
\end{tikzpicture}
}
\caption{Entity-relationship diagram of the main tables}
\label{fig:er}
\end{figure}

Figure \ref{fig:er} shows the main entities, created by Alembic migrations 0001 to 0012. The \texttt{users} table holds the account, the basic financial profile, the deletion schedule and the AI consent flag. \texttt{financial\_profiles} has a one-to-one relationship with users and stores gold, deposits, savings, property, loans, rent received and the monthly SIP, while \texttt{portfolios} holds shares and mutual funds. Each analysis is one row of \texttt{decisions}, which stores both the inputs (including snapshots of the portfolio and the comment reading as JSONB) and every model output. \texttt{audit\_logs} holds the hash-chained entries, and its foreign keys use \texttt{ON DELETE SET NULL} so that the chain survives when an account is anonymized. Row-level security is enabled and forced on decisions, audit logs, portfolios, financial profiles and the profile change log.

\section{Deployment Diagram}
\begin{figure}[H]
\centering
\resizebox{0.9\textwidth}{!}{\begin{tikzpicture}
  \node[box, text width=40mm, minimum height=16mm] (vite) {\textbf{Vite dev server :3000}\\ \footnotesize React SPA, proxies /v1};
  \node[box, right=8mm of vite, text width=40mm, minimum height=16mm] (uv) {\textbf{Uvicorn :8000}\\ \footnotesize FastAPI app and\\ \footnotesize background pipeline};
  \node[mlbox, right=8mm of uv, text width=34mm, minimum height=16mm] (art) {\textbf{Files on disk}\\ \footnotesize models/, results/, data/};
  \node[store, below=12mm of uv.south west, anchor=north east, xshift=-2mm, text width=40mm] (pg) {\textbf{PostgreSQL 18 :5432}\\ \footnotesize finai (owner), app\_user (RLS)};
  \node[store, below=12mm of uv.south east, anchor=north west, xshift=2mm, text width=40mm] (rd) {\textbf{Memurai / Redis :6379}};
  \draw[biflow] (vite) -- (uv);
  \draw[biflow] (uv) -- (art);
  \draw[biflow] (uv.south) ++(-5mm,0) -- (pg.north);
  \draw[biflow] (uv.south) ++(5mm,0) -- (rd.north);
  \begin{scope}[on background layer]
    \node[tier, fit=(vite)(art)(pg)(rd), inner sep=5mm] (dev) {};
  \end{scope}
  \node[tierlabel, anchor=south west] at (dev.north west) {Device: development and demonstration laptop (Windows, Python 3.13)};
\end{tikzpicture}
}
\caption{Deployment diagram of the development set-up}
\label{fig:deploy}
\end{figure}

Figure \ref{fig:deploy} shows the configuration used for development and demonstration, with all processes on one laptop. PostgreSQL has two roles: \texttt{finai}, the owner used for migrations, and \texttt{app\_user}, a non-superuser used by the application so that the row-level security policies apply. Container definitions are present in the repository for the planned public deployment.

\section{User Interface Design}
The public pages are the landing page, the privacy policy and the sign-in pages, together with \texttt{/auth/google}, where Google sign-in returns. A first Google sign-in lands on a Welcome page that asks for the basic profile; until it is filled, every other signed-in page except Settings sends the investor there. The application pages are an overview dashboard, the portfolio, a new-analysis form with a live market card, the result page, the history and the settings.

The result page is the centre of the design. It shows the live five-step progress with plain titles such as ``Reading your decision'' and ``Projecting the future'', then the recommendation with its reasons, the position against the whole portfolio, a chart of both futures, the pattern found with probability bars, the six-part explanation with an English and Hindi switch, report downloads, the inputs that weighed most, cause and effect, both sides in numbers, the investor's own words, ``Verify this record'' and ``How this was worked out''. The interface never shows a model name or an accuracy figure, and confidence is shown in words.
